\documentclass[10pt,a4paper]{amsart}
\usepackage[margin=19mm]{geometry}
\usepackage{amsmath,amssymb,mathtools}
\usepackage[scr=rsfs,cal=euler]{mathalfa}
\usepackage{xfrac}
\usepackage[unicode,hidelinks,bookmarks=false]{hyperref}
\newcommand{\ie}{i.e.\ }
\newcommand{\cf}{cf.\ }
\newcommand{\resp}{respectively\ }
\newcommand{\Subsection}{\subsection}
\providecommand{\nobreakdash}{\nobreak\hbox{-}}
\setlength{\emergencystretch}{2em}
\multlinegap0pt
\usepackage{bm}
\usepackage{bbm}
\usepackage{dsfont}
\usepackage{xspace}
\usepackage{cancel}
\usepackage{mathtools}
\usepackage{amsthm}
\makeatletter
\@ifundefined{equation}{\newtheorem{equation}{Equation}}{}
\@ifundefined{lemma}{\newtheorem{lemma}[equation]{Lemma}}{}
\@ifundefined{theorem}{\newtheorem{theorem}[equation]{Theorem}}{}
\makeatother
\title[Variable exponent modulus in symmetric domains]{Variable exponent modulus in symmetric domains}
\author{Rahim Kargar}
\date{Source: arXiv:2603.26941v2 (CC BY 4.0)}
\begin{document}
\maketitle

\noindent\emph{Source and licence.} Variable exponent modulus in symmetric domains. Version arXiv:2603.26941v2 (CC BY 4.0), distributed under the Creative Commons Attribution 4.0 International licence, as recorded in the arXiv licence metadata for this exact version and shown on the abstract page https://arxiv.org/abs/2603.26941v2. Full rights notice: licenses/arxiv-2603.26941-CC-BY-4.0.txt.\par\smallskip

\noindent\textit{Editorial scope.} Counted unit: the Theorem whose source label is \texttt{thm:annulus\_formula} in the article (source0.tex lines 1704--1750, source label \texttt{thm:annulus\_formula}), delivered together with the article's own complete original proof of it (source0.tex lines 1752--1945, one \texttt{proof} environment of 194 lines). No mathematics is written, repaired or re-derived: statement and proof are the article's own text line for line. Required context retained uncounted: fun-h (lines 937--943); lem:radial\_suffices (lines 1242--1307); thm:1d\_modulus (lines 1474--1505). Every definition, display and label of those lines is retained unchanged. Omitted from this package and not redistributed: 1--936, 944--1241, 1308--1473, 1506--1703, 1751--1751, 1946--3165. Nothing inside a retained excerpt is omitted or altered: every retained body is delivered exactly as the article's lines are written, so no omission receipt of any kind accompanies this package. Citations: the retained excerpts carry no citation command at all, so no bibliography entry of the article is transcribed and no reference is redistributed. Reference adaptation: none was needed and none was made; every reference inside the delivered unit resolves inside it, so no unresolved reference or citation is produced. Printed slips kept literal: no printed word, symbol, hypothesis or number of the counted statement or of its proof was altered. Layout and class: the article's own class is \texttt{amsart} with packages including amsmath, amssymb, amsthm, mathrsfs, babel, fontenc, graphicx, a4wide, float, hyperref; the wrapper is the lane's pinned \texttt{amsart} editorial wrapper at 10pt on A4, which restates the theorem-like environments the retained text needs and adds the article's own macro definitions that the retained text needs (none), with extra wrapper packages (bm, bbm, dsfont, xspace, cancel, mathtools). The delivered numbering is the unit's own. Assets: no retained span contains a figure environment, a graphics-inclusion command, a PDF-page inclusion, a table or a TikZ picture; no image file is copied into this package, the package declares no asset dependency and no image file is redistributed with it. Licence: version arXiv:2603.26941v2 of this article is distributed under the Creative Commons Attribution 4.0 International licence, as recorded in the arXiv licence metadata for this exact version and shown on the abstract page https://arxiv.org/abs/2603.26941v2. A full rights notice is delivered at licenses/arxiv-2603.26941-CC-BY-4.0.txt. Characters: every retained excerpt is pure ASCII, so the delivered engine, XeLaTeX with its default Latin Modern text font, reports no missing character.
\begin{equation}\label{fun-h}
h
=
\frac{\chi_S}{|S|}
-
\frac{\chi_T}{|T|}.
\end{equation}

\begin{lemma}
\label{lem:radial_suffices}
Let
\begin{equation*}
A(r_1,r_2)
=
\{x\in\mathbb R^n:r_1<|x|<r_2\},
\qquad
0<r_1<r_2<\infty,
\end{equation*}
and assume that
$p(x)=q_0(|x|)$
for a measurable function
$q_0:(r_1,r_2)\longrightarrow(1,\infty)$
satisfying
\begin{equation}
\label{eq:annulus_exp_bounds}
1<q_0^-\le q_0(r)\le q_0^+<\infty
\end{equation}
for almost every $r\in(r_1,r_2)$.
Let $\Gamma$ be the family of locally rectifiable curves in
$\overline{A(r_1,r_2)}$ whose endpoints belong to
$S_{r_1}^{n-1}$ and $S_{r_2}^{n-1}$, and whose interiors
lie in $A(r_1,r_2)$.

For $\rho\in\mathcal F(\Gamma)$, define its spherical average by
\begin{equation}
\label{eq:spherical_average}
\varphi(r)
=
\frac{1}{\omega_{n-1}}
\int_{S^{n-1}}
\rho(r\theta)\,d\sigma(\theta),
\qquad
r\in(r_1,r_2),
\end{equation}
and set
\begin{equation}
\label{eq:radial_average_density}
\widetilde\rho(x)
=
\varphi(|x|).
\end{equation}
Then $\varphi$ is Lebesgue measurable,
$\widetilde\rho\in\mathcal F(\Gamma)$, and
\begin{equation}
\label{eq:radial_energy_reduction}
\int_{A(r_1,r_2)}
\widetilde\rho(x)^{p(x)}\,dx
\le
\int_{A(r_1,r_2)}
\rho(x)^{p(x)}\,dx.
\end{equation}
Consequently,
\begin{equation}
\label{eq:radial_infimum}
M_{p(\cdot)}(\Gamma)
=
\inf\left\{
\int_{A(r_1,r_2)}
\rho(x)^{p(x)}\,dx:
\rho\in\mathcal F(\Gamma),
\ \rho\text{ is radial}
\right\}.
\end{equation}
\end{lemma}

\begin{theorem}
\label{thm:1d_modulus}
Under the assumptions of Lemma~\ref{lem:radial_suffices}, the modular
variable exponent modulus admits the one-dimensional representation
\begin{equation}
\label{eq:annulus_1d_reduction}
M_{p(\cdot)}(\Gamma)
=
\inf_{\varphi\in\mathcal A_r} J_r(\varphi),
\end{equation}
where
\begin{equation}
\label{eq:annulus_admissible_class}
\mathcal A_r
=
\left\{
\varphi\in L^{q_0(\cdot)}(r_1,r_2):
\ \varphi\ge0\ \text{a.e.},\
\int_{r_1}^{r_2}\varphi(r)\,dr=1
\right\},
\end{equation}
and
\begin{equation}
\label{eq:annulus_1d_energy}
J_r(\varphi)
=
\omega_{n-1}
\int_{r_1}^{r_2}
\varphi(r)^{q_0(r)}r^{n-1}\,dr.
\end{equation}
Moreover, $J_r$ admits a unique minimizer in $\mathcal A_r$.
\end{theorem}

\begin{theorem}
\label{thm:annulus_formula}
Under the assumptions of Lemma~\ref{lem:radial_suffices}, let
$\rho_*$ denote the unique minimizer of $J_r$ in $\mathcal A_r$.
Then $\rho_*>0$ almost everywhere in $(r_1,r_2)$ and there exists a
unique constant $\lambda>0$ such that
\begin{equation}
\label{eq:annulus_EL}
\omega_{n-1}q_0(r)r^{n-1}
\rho_*(r)^{q_0(r)-1}
=
\lambda
\end{equation}
for almost every $r\in(r_1,r_2)$.
Consequently,
\begin{equation}
\label{eq:annulus_extr}
\rho_*(r)
=
\left(
\frac{\lambda}
{\omega_{n-1}q_0(r)r^{n-1}}
\right)^{\!1/(q_0(r)-1)}
\end{equation}
for almost every $r\in(r_1,r_2)$,
where $\lambda>0$ is uniquely determined by
\begin{equation}
\label{eq:annulus_normalization}
\int_{r_1}^{r_2}
\left(
\frac{\lambda}
{\omega_{n-1}q_0(r)r^{n-1}}
\right)^{\!1/(q_0(r)-1)}
\,dr
=
1.
\end{equation}
Moreover,
\begin{equation}
\label{eq:annulus_modulus_formula}
M_{p(\cdot)}(\Gamma)
=
\omega_{n-1}
\int_{r_1}^{r_2}
\rho_*(r)^{q_0(r)}r^{n-1}\,dr.
\end{equation}
\end{theorem}

\begin{proof}
Let $\rho_*$ be the unique minimizer of $J_r$ in $\mathcal A_r$, whose
existence and uniqueness were established in
Theorem~\ref{thm:1d_modulus}. We first note that
\begin{equation}
\label{eq:annulus_positive}
\rho_*(r)>0
\end{equation}
for almost every $r\in(r_1,r_2)$.
Indeed, the positivity argument used for the unweighted one-dimensional
problem applies without change to the present functional, because
\begin{equation*}
0<\omega_{n-1}r_1^{n-1}
\le
\omega_{n-1}r^{n-1}
\le
\omega_{n-1}r_2^{n-1}
<\infty
\end{equation*}
on $(r_1,r_2)$. If $\rho_*$ vanished on a set of positive measure, one
could transfer a sufficiently small amount of mass from a subset on
which $\rho_*$ is bounded below to that zero set. The contribution of
the zero set would be of order $\varepsilon^{q_0^-}$, whereas the
decrease on the positive set would be of order $\varepsilon$. Since
$q_0^->1$, this would strictly decrease $J_r$, contradicting minimality.

We next establish the Euler--Lagrange equation. Define
\begin{equation*}
g(r)
:=
\omega_{n-1}q_0(r)r^{n-1}
\rho_*(r)^{q_0(r)-1}.
\end{equation*}
Fix $0<\delta<N<\infty$ and consider the measurable level set
\begin{equation*}
E_{\delta,N}
:=
\{r\in(r_1,r_2):\delta\le\rho_*(r)\le N\}.
\end{equation*}
By \eqref{eq:annulus_positive}, these sets exhaust $(r_1,r_2)$ up to a
null set as $\delta\downarrow0$ and $N\uparrow\infty$.

Let $S,T\subset E_{\delta,N}$ be disjoint measurable sets of positive
measure. Define $h$ as in \eqref{fun-h}.
Then $\int_{r_1}^{r_2}h(r)\,dr=0.$
For sufficiently small $|\varepsilon|$, the function
$\rho_\varepsilon=\rho_*+\varepsilon h$
is nonnegative and satisfies
$\int_{r_1}^{r_2}\rho_\varepsilon(r)\,dr=1.$
Hence, $\rho_\varepsilon\in\mathcal A_r$ for both signs of sufficiently
small $\varepsilon$.

On $S\cup T$, the functions $\rho_\varepsilon$ remain in a fixed
compact subinterval of $(0,\infty)$ for sufficiently small
$|\varepsilon|$. Since $q_0$ is uniformly bounded between $q_0^-$ and
$q_0^+$, the difference quotients
\begin{equation*}
\frac{
1
}{\varepsilon}\left(\rho_\varepsilon(r)^{q_0(r)}
-
\rho_*(r)^{q_0(r)}\right)
\end{equation*}
are dominated by an integrable function independent of $\varepsilon$.
Therefore, differentiation under the integral sign is justified, and
minimality of $\rho_*$ for both positive and negative $\varepsilon$
gives
\begin{equation}
\label{eq:annulus_first_variation}
0
=
\left.
\frac{d}{d\varepsilon}
J_r(\rho_\varepsilon)
\right|_{\varepsilon=0}
=
\int_{r_1}^{r_2}
g(r)h(r)\,dr.
\end{equation}
Consequently,
\begin{equation*}
\frac{1}{|S|}\int_S g(r)\,dr
=
\frac{1}{|T|}\int_T g(r)\,dr.
\end{equation*}
Since this holds for every pair of disjoint positive-measure subsets
$S,T\subset E_{\delta,N}$, it follows that $g$ is constant almost
everywhere on $E_{\delta,N}$. Thus, there exists a constant
$\lambda_{\delta,N}$ such that $g(r)=\lambda_{\delta,N}$
for almost every $r\in E_{\delta,N}$.

If two such level sets have an intersection of positive measure, their
corresponding constants must coincide. Since the sets
$E_{\delta,N}$ exhaust $(r_1,r_2)$ up to a null set, there is therefore
a single constant $\lambda\in\mathbb R$ such that $g(r)=\lambda$ for almost every $r\in(r_1,r_2)$.
This proves \eqref{eq:annulus_EL}.

In view of
\eqref{eq:annulus_positive}, together with
$q_0(r)>1$ and $r>0$, every factor on the left-hand side of
\eqref{eq:annulus_EL} is positive almost everywhere. Hence
$\lambda>0$.
Solving \eqref{eq:annulus_EL} for $\rho_*(r)$ gives
\eqref{eq:annulus_extr}. Since $\rho_*\in\mathcal A_r$, its normalization
immediately yields \eqref{eq:annulus_normalization}.

It remains to prove that the normalization equation determines $\lambda$
uniquely. Define
\begin{equation}
\label{eq:annulus_G}
G(\lambda)
:=
\int_{r_1}^{r_2}
\left(
\frac{\lambda}
{\omega_{n-1}q_0(r)r^{n-1}}
\right)^{\!1/(q_0(r)-1)}
\,dr,
\qquad \lambda>0.
\end{equation}
For each fixed $r$, the integrand is continuous and strictly increasing
as a function of $\lambda>0$. Hence $G$ is strictly increasing.

We next verify continuity. Put
\begin{equation*}
d_-
:=
\omega_{n-1}q_0^-r_1^{n-1},
\qquad
d_+
:=
\omega_{n-1}q_0^+r_2^{n-1}.
\end{equation*}
More directly, since
\begin{equation*}
\omega_{n-1}q_0^-r_1^{n-1}
\le
\omega_{n-1}q_0(r)r^{n-1}
\le
\omega_{n-1}q_0^+r_2^{n-1},
\end{equation*}
the quotient appearing in \eqref{eq:annulus_G} is uniformly bounded
above on every compact interval $[\lambda_1,\lambda_2]\subset(0,\infty)$.
Indeed, if
\begin{equation*}
a_1=\frac{\lambda_1}
{\omega_{n-1}q_0^+r_2^{n-1}},
\qquad
a_2=\frac{\lambda_2}
{\omega_{n-1}q_0^-r_1^{n-1}},
\end{equation*}
then the quotient belongs to $[a_1,a_2]$, while
\begin{equation*}
\frac{1}{q_0^+-1}
\le
\frac{1}{q_0(r)-1}
\le
\frac{1}{q_0^--1}.
\end{equation*}
Thus, the integrand is bounded by a finite constant depending only on
$\lambda_2$, $q_0^\pm$, $r_1$, $r_2$, and $\omega_{n-1}$. Since
$(r_1,r_2)$ has finite measure, dominated convergence proves that
$G$ is continuous on $(0,\infty)$.

As $\lambda\downarrow0$, the integrand in \eqref{eq:annulus_G}
converges pointwise to zero. For $0<\lambda\le1$, it is bounded above by
the integrand corresponding to $\lambda=1$, which is integrable by the
same uniform bounds. Hence
$\lim_{\lambda\downarrow0}G(\lambda)=0.$
On the other hand, as $\lambda\uparrow\infty$, the integrand converges
pointwise to $+\infty$. By the monotone convergence theorem,
$\lim_{\lambda\uparrow\infty}G(\lambda)=+\infty.$
Therefore, the intermediate value theorem and strict monotonicity imply
that there exists a unique $\lambda>0$ satisfying $G(\lambda)=1$.
This proves the uniqueness asserted in \eqref{eq:annulus_normalization}.

Finally, by Theorem~\ref{thm:1d_modulus},
\begin{equation*}
M_{p(\cdot)}(\Gamma)
=
\inf_{\varphi\in\mathcal A_r}J_r(\varphi)
=
J_r(\rho_*),
\end{equation*}
and therefore
\begin{equation*}
M_{p(\cdot)}(\Gamma)
=
\omega_{n-1}
\int_{r_1}^{r_2}
\rho_*(r)^{q_0(r)}r^{n-1}\,dr.
\end{equation*}
This proves \eqref{eq:annulus_modulus_formula} and we are done.
\end{proof}

\end{document}
